% !TeX root = ../../../main.tex
\subsection{Log-Time Cumulative Distribution Function}
\label{ch06:apd:ltcdf}


Recall the log-time cumulative distribution function is given as follows:

\[
   F(x) = \frac{1}{k} \sum_{i=1}^k \mathbf{1}_{\{x_{(i)} \leq x\}}
\]
where $\mathbf{1}$ is the indicator function.



We compute the AULTC as the area under $F(x)$ from $x = 0$ to $x = \log(1 + S_{\text{max}})$, normalized by $\log(1 + S_{\text{max}})$:
\begin{align*}
\text{AULTC} &= \frac{1}{\log(1 + S_{\text{max}})} \Bigg[ \sum_{i=1}^k (x_{(i)} - x_{(i-1)}) \frac{i}{k} + (\log(1 + S_{\text{max}}) - x_{(k)}) \times 1 \Bigg]
\end{align*}
where $x_{(0)} = 0$.
With this definition, AULTC = 1 indicates that discrepancies are zero (perfect recovery), resulting in maximum area $\log(1 + S_{\text{max}})$, and AULTC = 0 indicates that all discrepancies exceed $S_{\text{max}}$, yielding zero area. 
The code for the AULTC calculation is given in the function \texttt{ecdf\textunderscore auc2(\dots)} in the Supplementary Materials.


\begin{ChSixRemark}
The time unit and cutoff $S_\text{max}$ affect the AULTC calculation, so they must be specified when reporting the AULTC.
\end{ChSixRemark}
Unless otherwise stated, we set $S_{\max}=8760$ hours, corresponding to one year.
This upper limit was chosen because it captures clinically meaningful discrepancies from hours to months while preventing very long-horizon case-history events from dominating the metric.

Because of the $\log(1 + \cdot)$ transformation (rather than the $\log(\cdot)$ transformation which is undefined for zero time-error discrepancy), the discrepancies are non-linearly shifted in the log scale.
The non-linear shift adjusts the relative widths of the step function (particularly for small discrepancies) changing the area calculation.
The cutoff $S_\text{max}$ affects the normalization factor.
Therefore, these values should be chosen based on practicalities for the application.
In our case, we chose hours at the time unit because the only sub-hour descriptions reported were several Apgar scores at the minute-level.

\begin{ChSixRemark}
    The average log-time discrepancy is non-convex.
\end{ChSixRemark}
This can be seen from the observation that the $\log(1+\cdot)$ function is ``spikier'' than the L1 function.
More formally, the log-time discrepancy has a higher curvature around zero error compared to the L1 loss.
Consider the second derivatives of the loss functions with respect to the discrepancies $s = |t^p - t^r|$.
For L1 loss $L_1(s) = s$, the second derivative $\frac{d^2 L_1}{ds^2} = 0$ for $s > 0$.
For the log-time loss, $L_{log}(s) = \log(1 + s)$, the second derivative is $\frac{d^2 L_{log}}{ds^2} = -\frac{1}{(1 + s)^2}$.
At $s = 0$, the second derivative is $-1$, while the second derivative of the L1 loss is $0$.

To provide a simple counterexample, define the average log-time discrepancy is given by:
\[
L(t^p) = \frac{1}{k} \sum_{i=1}^k \log(1 + \min(|t^p_i - t^r_i|, S_{\text{max}}))
\]
To prove that this function is non-convex, we need to show that there exist $t^{p1}, t^{p2}$ and $\lambda \in (0,1)$ such that
\[
L(\lambda t^{p1} + (1-\lambda) t^{p2}) > \lambda L(t^{p1}) + (1-\lambda) L(t^{p2})
\]
Consider the case where $k=1$, $t^r = 0$, and $S_{\text{max}} = 2$.
The loss function is $L(t^p) = \log(1 + \min(|t^p|, 2))$.
Let $t^{p1} = 0$ and $t^{p2} = 3$.
Choose $\lambda = 0.5$.
Then $\lambda t^{p1} + (1-\lambda) t^{p2} = 0.5 \times 0 + 0.5 \times 3 = 1.5$.
Now we evaluate the loss function at these points:
\begin{align*}
L(t^{p1}) &= L(0) = \log(1 + \min(|0|, 2)) = 0 \\
L(t^{p2}) &= L(3) = \log(1 + \min(|3|, 2)) = \\
    & \qquad \log(3) \\
L(\lambda t^{p1} + (1-\lambda) t^{p2}) &= L(1.5) = \log(2.5)
\end{align*}
Now we check the convexity condition:
\[
L(1.5) \leq 0.5 L(0) + 0.5 L(3)
\]
\[
\log(2.5) \leq 0.5 \times 0 + 0.5 \times \log(3)
\]
\[
\log(2.5) \leq \log(\sqrt{3})
\]
Since $2.5 > \sqrt{3} \approx 1.732$, we have $\log(2.5) > \log(\sqrt{3})$.
Therefore,
\[
L(\lambda t^{p1} + (1-\lambda) t^{p2}) > \lambda L(t^{p1}) + (1-\lambda) L(t^{p2})
\]
This proves that the average log-time discrepancy is non-convex.
This remark illustrates the loss form, in case it is considered for optimization/model training rather than for assessment purposes.










